\section{Introduction}
Outlier analysis is a critical task in both academic research and industrial applications. 
In recent years, representation learning has been increasingly adopted to automate and refine outlier detection. 
A notable contribution to this field is the InterCont method, an approach based on internal contrastive learning proposed by \citet{shenkar2022anomaly}. 
While outlier detection is fundamentally a classification task, it is characterized by extreme class imbalance.
Despite its importance, a single formal definition of an outlier has not reached universal consensus.
\citet{ayadi2017outlier} identified twelve distinct interpretations. 
This study adopts the definition established by \citet{hawkins1980identification}, which identifies outliers, also termed anomalies, \todo{Chandola et al., 2009 what are you about?}
as observations that deviate so significantly from other observations as to arouse suspicions that they were generated by a different mechanism. 
Depending on the research context, outliers may either represent operative failures (recording errors, misreporting, sampling errors) to be removed 
or primary targets of interest that provide critical information \citet{smiti2020critical}.


For example in cybersecurity, it facilitates the detection of network intrusions or malware through unusual traffic patterns. 
In finance, these methods identify fraudulent activities, such as the unauthorized use of credit cards \citep{rezapour2019anomaly}. 
In healthcare, outlier detection assists in diagnosing rare diseases and identifying anomalous patient data.
Taxonomically, outliers are categorized into three primary types:


\begin{itemize}
    \item   Global Outliers: Individual data points that deviate significantly from the entire dataset.
    \item   Contextual Outliers: Data points considered anomalous only within a specific context. 
    \item   Collective Outliers: A sequence of data points that are not anomalous individually but become so when appearing together.
\end{itemize}           


\subsubsection*{Delimitation}
Distinct from outlier detection are related tasks such as Noisy Data Identification, which addresses incorrect or missing values.
Novelty Detection, which seeks previously unobserved but potentially valid patterns \citet{domingues2019probabilistic}.
Out-of-Distribution Detection (OOD), which identifies data points that do not belong to any of the known training classes. \todo{more specific?}


A contemporary shift in the field emphasizes Explainable Outlier Detection. 
It is no longer sufficient to flag an anomaly.
Models must provide the underlying rationale for the classification. 
This often involves assigning importance levels based on the distance from the normal distribution or identifying causal interactions where one primary outlier triggers a sequence of secondary anomalies \citet{panjei2022survey}.
\todo{explain that ICL is explainable too but not the theme of this thesis}


\subsubsection*{Challenges}
High-dimensional datasets introduce the Curse of Dimensionality (CoD), which significantly complicates detection tasks. 
As dimensionality ($p$) increases, the volume of the space grows exponentially, causing data points to become sparsely distributed. 
This sparsity increases the mean distance between points, making the Euclidean distance—a common measure of similarity—less reliable. 
Furthermore, in cases where $p > n$, perfect multicollinearity may occur, where variables become linear combinations of others, leading to redundancy. 
High dimensionality also increases the risk of overfitting, where models capture stochastic noise rather than a generalizable signal \citet{altman2018curse}.
\todo{concept drift?}
\todo{scarcity of lables?}


A major challenge in outlier detection research involves establishing comparability between diverse methods, which necessitates the implementation of a unified framework. 
Such a framework ensures that the same datasets, paradigms, and evaluation metrics are consistently applied.
Utilizing benchmark datasets recognized by the scientific community, such as those provided by \citet{Rayana:2016}, facilitates this comparability across different methodologies. 
By applying algorithms to a comprehensive collection of datasets with varying properties, the resulting performance metrics reveal the specific advantages and disadvantages of each approach. 
In this thesis, the "Wine" dataset from the ODDS Repository is utilized, which implies that the findings are specifically constrained to the characteristics of this data. 




The exceptional performance of the $k$-nearest neighbors algorithm in the analysis of \citet{shenkar2022anomaly}, which surpasses the more complex architecture of InterCont and F1-score, serves as a primary catalyst for this research. 
The fact that a traditional, distance-based approach can yield superior results compared to sophisticated representation learning methods like InterCont raises critical questions about the necessity of algorithmic complexity in outlier detection. 
This empirical discrepancy between theoretical sophistication and practical efficacy directly inspired the formulation of the following hypotheses:


\begin{enumerate}
    \item H1: k-NN will achieve the highest predictive performance across the selected metrics.
    \item H2: InterCont will demonstrate superior performance relative to the statistical benchmark (MCD) but remain secondary to k-NN.
\end{enumerate}


This thesis will proceed with the following: 
The first chapter establishes the theoretical foundations of machine learning, beginning with a brief description of neural network architecture before transitioning to the advanced topic of contrastive learning. 
This section outlines the primary framework of contrastive learning through a computer vision example. 
The second chapter details the central methodology: Internal Contrastive Learning. 
The focus remains on the architecture of this method and the mechanisms through which it identifies outliers, concluding with an analysis of its advantages and disadvantages. 
The subsequent chapter introduces the first benchmark method, k-nearest neighbors (k-NN). 
This section explains the underlying mechanics and architecture of k-NN while justifying its selection as a benchmark. 
In the following chapter, the Minimum Covariance Determinant (MCD) is introduced. 
The discussion details how this method detects outliers and provides the rationale for its use as a benchmark against Internal Contrastive Learning.


In the data analysis chapter, the dataset is initially examined by describing its primary characteristics, which serve as the basis for interpreting and explaining the results. 
The analysis begins with a univariate approach, utilizing boxplots to describe parameters such as the mean and distribution. 
Subsequently, the analysis proceeds to a bivariate approach, where the data is described using a correlation heatmap. 
The final stage of the analysis involves characterizing the geometry of the dataset through UMAP, a dimension reduction technique. 
In the metrics chapter, the most significant metrics used in outlier detection research—facilitated by the confusion matrix—are described alongside a discussion of their respective advantages and disadvantages. 
The selection of metrics is finalized with the F1-Score, AUC ROC, and AUPRC. The results section presents the findings, followed by the discussion section, where the results are explained and interpreted. 
The limitations section provides a critical evaluation of the methodology and identifies specific constraints. 
The thesis concludes with a summary of the methodology and the principal findings.